\documentclass[10pt,a4paper]{amsart}
\usepackage{amsmath,amssymb,amsthm,graphicx}
% Retained minimal preamble for the Stacks Project chapter more-morphisms.tex (More on Morphisms).
% Verbatim excerpt of the chapter preamble: Xy-pic setup (native preamble lines 14-18) and the native
% theorem declarations with the mathematical macros (native preamble lines 154-217). Copied mechanically;
% nothing added, removed or redefined. The wrapper class and the amsmath/amssymb/amsthm/graphicx base
% packages are supplied by the assembler, so the class/externaldocument part of the native preamble is omitted.
% For commutative diagrams we use Xy-pic
\usepackage[all]{xy}

% We use 2cell for 2-commutative diagrams.
\xyoption{2cell}
% Theorem environments.
%
\theoremstyle{plain}
\newtheorem{theorem}[subsection]{Theorem}
\newtheorem{proposition}[subsection]{Proposition}
\newtheorem{lemma}[subsection]{Lemma}

\theoremstyle{definition}
\newtheorem{definition}[subsection]{Definition}
\newtheorem{example}[subsection]{Example}
\newtheorem{exercise}[subsection]{Exercise}
\newtheorem{situation}[subsection]{Situation}

\theoremstyle{remark}
\newtheorem{remark}[subsection]{Remark}
\newtheorem{remarks}[subsection]{Remarks}

\numberwithin{equation}{subsection}

% Macros
%
\def\lim{\mathop{\mathrm{lim}}\nolimits}
\def\colim{\mathop{\mathrm{colim}}\nolimits}
\def\Spec{\mathop{\mathrm{Spec}}}
\def\Hom{\mathop{\mathrm{Hom}}\nolimits}
\def\Ext{\mathop{\mathrm{Ext}}\nolimits}
\def\SheafHom{\mathop{\mathcal{H}\!\mathit{om}}\nolimits}
\def\SheafExt{\mathop{\mathcal{E}\!\mathit{xt}}\nolimits}
\def\Sch{\mathit{Sch}}
\def\Mor{\mathop{\mathrm{Mor}}\nolimits}
\def\Ob{\mathop{\mathrm{Ob}}\nolimits}
\def\Sh{\mathop{\mathit{Sh}}\nolimits}
\def\NL{\mathop{N\!L}\nolimits}
\def\CH{\mathop{\mathrm{CH}}\nolimits}
\def\proetale{{pro\text{-}\acute{e}tale}}
\def\etale{{\acute{e}tale}}
\def\QCoh{\mathit{QCoh}}
\def\Ker{\mathop{\mathrm{Ker}}}
\def\Im{\mathop{\mathrm{Im}}}
\def\Coker{\mathop{\mathrm{Coker}}}
\def\Coim{\mathop{\mathrm{Coim}}}

% Boxtimes
%
\DeclareMathSymbol{\boxtimes}{\mathbin}{AMSa}{"02}

%
% Macros for moduli stacks/spaces
%
\def\QCohstack{\mathcal{QC}\!\mathit{oh}}
\def\Cohstack{\mathcal{C}\!\mathit{oh}}
\def\Spacesstack{\mathcal{S}\!\mathit{paces}}
\def\Quotfunctor{\mathrm{Quot}}
\def\Hilbfunctor{\mathrm{Hilb}}
\def\Curvesstack{\mathcal{C}\!\mathit{urves}}
\def\Polarizedstack{\mathcal{P}\!\mathit{olarized}}
\def\Complexesstack{\mathcal{C}\!\mathit{omplexes}}
% \Pic is the operator that assigns to X its picard group, usage \Pic(X)
% \Picardstack_{X/B} denotes the Picard stack of X over B
% \Picardfunctor_{X/B} denotes the Picard functor of X over B
\def\Pic{\mathop{\mathrm{Pic}}\nolimits}
\def\Picardstack{\mathcal{P}\!\mathit{ic}}
\def\Picardfunctor{\mathrm{Pic}}
\def\Deformationcategory{\mathcal{D}\!\mathit{ef}}
\newenvironment{reference}{}{}
\newenvironment{slogan}{}{}
\newenvironment{history}{}{}
\usepackage{cleveref}
\crefname{originallemma}{original lemma}{original lemmas}
\Crefname{originallemma}{Original lemma}{Original lemmas}
\makeatletter
\newlabel{algebra-lemma-unique-prime-over-localize-below}{{00EA}{source}{}{sourcepointer-1}{}}
\newlabel{algebra-lemma-unique-prime-over-localize-below@cref}{{[originallemma][0][]00EA}{[1][0][]source}{}{}{}}
\makeatother
\title{A finite morphism with a singleton fibre is locally determined near the point, with finite local rings and matching stalks}
\author{The Stacks Project Authors}
\date{stacks-project:more-morphisms.tex:local-snapshot-20261011; GNU FDL 1.2 or later}
\begin{document}
\maketitle
\noindent\textit{Editorial note.} Statement span is lines 13551-13568 of sources/stacks-more-morphisms/pinned-original-source.tex; proof span is lines 13570-13591, the complete author proof (20 body lines).\par\medskip
\noindent\textit{Editorial note.} Navigation stubs: 1 cited-but-not-included labels are supplied as non-load-bearing pointers with their pinned tags where the snapshot tag index has them.\par\medskip
\noindent\textit{Editorial note.} extra\_packages is empty on purpose: the retained preamble already loads xy with options and a second declaration clashes in the assembler wrapper.\par\medskip
\noindent\textit{Editorial note.} Text is retained verbatim, including the author's own typos and placeholders.\par\medskip
\noindent\textit{Editorial note.} Workhorse lemma used by the chapter's etale-local structure arguments; its parenthetical (some details omitted) covers only routine localization properties while every claim is named.\par\medskip
\noindent\textit{Allowed standard prerequisites (editorial).}\begin{itemize}
\item Explicit prerequisite boundary. Basic properties of finite morphisms (closed, affine), the unique-prime-over localization lemma (Algebra, lemma-unique-prime-over-localize-below), and elementary localization identities for modules.
\end{itemize}\medskip
\ifdefined\hypertarget\hypertarget{sourcepointer-1}{}\fi
\noindent\textit{Source reference.} Original lemma 00EA: sources/stacks-more-morphisms/pinned-original-source.tex, source lines 13591--13591. Cited by the retained statement or proof but not part of the collected spans; supplied only as a navigation pointer. The argument does not depend on its content.\par\medskip
\begin{lemma}
\label{lemma-finite-morphism-single-point-in-fibre}
Let $\pi : X \to Y$ be a finite morphism.
Let $x \in X$ with $y = \pi(x)$ such that $\pi^{-1}(\{y\}) = \{x\}$.
Then
\begin{enumerate}
\item For every neighbourhood $U \subset X$ of $x$ in $X$, there
exists a neighbourhood $V \subset Y$ of $y$ such that
$\pi^{-1}(V) \subset U$.
\item The ring map $\mathcal{O}_{Y, y} \to \mathcal{O}_{X, x}$
is finite.
\item If $\pi$ is of finite presentation, then
$\mathcal{O}_{Y, y} \to \mathcal{O}_{X, x}$ is of finite presentation.
\item For any quasi-coherent $\mathcal{O}_X$-module $\mathcal{F}$
we have $\mathcal{F}_x = \pi_*\mathcal{F}_y$ as
$\mathcal{O}_{Y, y}$-modules.
\end{enumerate}
\end{lemma}

\begin{proof}
The first assertion is purely topological; use that
$\pi$ is a continuous and closed map such that $\pi^{-1}(\{y\}) = \{x\}$.
To prove the second and third parts we may assume
$X = \Spec(B)$ and $Y = \Spec(A)$. Then
$A \to B$ is a finite ring map and $y$ corresponds to a prime
$\mathfrak p$ of $A$ such that there exists a unique prime $\mathfrak q$ of
$B$ lying over $\mathfrak p$. Then
$B_{\mathfrak q} = B_{\mathfrak p}$, see
Algebra, Lemma \ref{algebra-lemma-unique-prime-over-localize-below}.
In other words, the map $A_{\mathfrak p} \to B_{\mathfrak q}$
is equal to the map $A_{\mathfrak p} \to B_{\mathfrak p}$ you get
from localizing $A \to B$ at $\mathfrak p$.
Thus (2) and (3) follow from simple properties of localization
(some details omitted). For the final statement, suppose that
$\mathcal{F} = \widetilde M$ for some $B$-module $M$.
Then $\mathcal{F} = M_{\mathfrak q}$ and
$\pi_*\mathcal{F}_y = M_{\mathfrak p}$. By the above these
localizations agree. Alternatively you can use part (1) and
the definition of stalks to see that $\mathcal{F}_x = \pi_*\mathcal{F}_y$
directly.
\end{proof}

\end{document}
